\section{Introduction}

\begin{figure}[b]
\centering
\begin{minipage}[t]{0.49\textwidth}
  \centering
  \includegraphics[width=\linewidth]{figures/scaling_instances.pdf}
\end{minipage}
\hfill
\begin{minipage}[t]{0.49\textwidth}
  \centering
  \includegraphics[width=\linewidth]{figures/scaling_performance.pdf}
\end{minipage}
\caption{
\textbf{Scaling task instances} (left) and \textbf{performance} (right) for SWE-agent's with \bugs{}.
Using \bugs{}, we can create $100$s to $1000$s of instances for any Python codebase, enabling us to train \texttt{SWE-agent-LM-32B} which achieves $40.2$\% on SWE-bench Verified.
}
\label{fig:scaling_teaser}
\end{figure}

Language Model (LM) agents, such as SWE-agent~\citep{yang_swe-agent_2024} or OpenHands~\citep{wang_openhands_2024}, 
have made remarkable progress towards automating software engineering (SE) tasks, as tracked by benchmarks such as SWE-bench~\citep{jimenez_swe-bench_2024}.
However, the most effective agents still rely on proprietary LMs, as building open source LMs for SE remains bottlenecked by the lack of large-scale, high-quality training data.
To ensure that open research remains relevant in this field, it is critical to develop infrastructure for collecting software engineering training data at scale.

The current open-source ecosystem offers two kinds of data sources to train LMs on SE tasks.
One simple approach is to crawl pull requests (PRs) and issues from GitHub repositories.
However, without execution environments or tests, these instances offer no reliable way of validating generated solutions,
and LMs are limited to learning from the surface form of code~\citep{xie2025swefixertrainingopensourcellms} or via rewards based on superficial string similarity~\citep{wei2025swerladvancingllmreasoning}.

In contrast, SWE-bench provides reliable validation by running unit tests against proposed solutions.
Another line of work has simply extended the SWE-bench collection strategy to a new set of repositories for training purposes \citep{pan_training_2024}. 
This produces flexible environments for training and distilling LM agents, since we can generate agent trajectories and filter them based on the unit test results.
% , resulting in the most promising open-source LMs for software engineering tasks to date.
However, the scalability of this approach is severely limited by the challenges associated with SWE-bench's collection strategy.
SWE-bench's filtering process leaves only a small number of PRs that not only resolve a Github issue, but also make meaningful changes to unit tests.
Also, setting up execution environments for each instance requires a substantial amount of human intervention.

In this paper, we introduce the \bugs{} toolkit, which marries the flexible execution environments of SWE-bench with scalable instance collection (Figure~\ref{fig:scaling_teaser}).
\bugs{} features several techniques to automatically synthesize bugs in existing GitHub repositories,
such as (1) generating errant rewrites of functions with an LM, (2) procedurally modifying the abstract syntax tree (AST) of functions, (3) undoing PRs, and (4) combining bugs.
Our key insight is that execution-based validation can not only validate proposed solutions,
but also identify bug candidates which cause substantial software regression (i.e., break tests).

\begin{figure}[t]
    \centering
    \includegraphics[width=\textwidth]{figures/preview.pdf}
    \caption{
    \bugs{} creates training data for software engineering agents by crafting bugs into real codebases.
    Given a codebase, we employ several strategies to create task instances that break existing tests.
    Using \bugs{}, we create $50$k+ task instances with execution environments from $128$ real world repositories.
    }
    \label{fig:preview}
\end{figure}

In a nutshell, \bugs{} puts forth the following task creation workflow, as shown in Figure~\ref{fig:preview}.
Given a codebase, we automatically set up a corresponding environment using SWE-agent~\citep{yang_swe-agent_2024}.
Within this environment, we then use the aforementioned techniques to synthesize $100$s to $1,000$s of task instances.
Finally, we craft realistic issue descriptions automatically with LMs.
\bugs{}'s design significantly reduces the amount of human labor and storage required for constructing execution environments.
% Beyond this, we alleviate the requirement to connect PRs with GitHub issues by crafting realistic issue descriptions automatically with LMs.
% Finally, we eliminate the human effort of setting up new repositories by demonstrating that this task can already be largely automated with existing software engineering agents \citep{yang_swe-agent_2024}, enabling scalable environment creation across many repositories.
% In addition, \bugs{} only requires a single execution environment for all of its instances, drastically reducing storage needs and further simplifying the workflow.
Using \bugs{}, we create a dataset of $50$k task instances across $128$ real-world GitHub repositories.

Using the \bugs{} dataset, we achieve a new open-weight state of the art result on SWE-bench verified.
Using the SWE-smith task instances, we generate $5{,}016$ expert trajectories with Claude 3.7 Sonnet and fine-tune Qwen 2.5 Coder Instruct $32$B.
The resulting LM,  {\tt SWE-agent-LM-32B}, achieves $40.2\%$ (+$33.4$\%) on SWE-bench Verified in a single attempt, without inference-time scaling.
This sets a new state of the art for open-weight models.

The scale and diversity of the \bugs{} dataset enables us to begin establishing truths and investigate interesting phenomena about developing SWE-agents.
Training on more instances, bug types, and repositories helps.
LM generated issue text approximates real ones effectively.
Using \bugs{}, we find that it's possible to optimize LMs to perform well for specific repositories while only suffering minor generalization loss.

We release \bugs{} as an open-source toolkit --- including instances, environments, and trajectories --- to catalyze the development of stronger open-source LM agents.
